\chapter{Đánh giá và Kết quả Thực nghiệm}

\begin{tcolorbox}[warnbox]
Toàn bộ kết quả trong chương này được thu thập trên \textbf{dữ liệu tổng hợp (synthetic mock data)} của phiên bản MVP --- mạng 20 nút giao, 1.000 aggregate producers, chu kỳ 5 phút trong 14 ngày. Các chỉ số chỉ mang tính tham chiếu; kết quả trên dữ liệu thực tế có thể khác biệt đáng kể về biên độ và phân bố.
\end{tcolorbox}

%% ============================================================
\section{Môi trường Thực nghiệm}
%% ============================================================

\subsection{Phần cứng}

Bảng~\ref{tab:hw_env} mô tả cấu hình phần cứng dùng cho toàn bộ benchmark.

\begin{table}[H]
\centering
\caption{Cấu hình phần cứng thực nghiệm}
\label{tab:hw_env}
\renewcommand{\arraystretch}{1.4}
\begin{tabular}{p{3.5cm} p{9cm}}
\toprule
\textbf{Thành phần} & \textbf{Chi tiết} \\
\midrule
CPU & Intel Core i7-12700H --- 14 cores / 20 threads, 2.3 GHz base \\
\rowcolor{rowgray}
RAM & 32 GB DDR5-4800 \\
GPU & NVIDIA RTX 3060 Laptop --- 6 GB VRAM (CUDA 12.2, cuDNN 8.9) \\
\rowcolor{rowgray}
Ổ cứng & NVMe SSD 512 GB (seq read $\sim$3.500 MB/s) \\
\bottomrule
\end{tabular}
\end{table}

\subsection{Phần mềm}

\begin{table}[H]
\centering
\caption{Ngăn xếp phần mềm thực nghiệm}
\label{tab:sw_env}
\renewcommand{\arraystretch}{1.4}
\begin{tabular}{p{3.5cm} p{9cm}}
\toprule
\textbf{Thành phần} & \textbf{Phiên bản} \\
\midrule
Hệ điều hành & Ubuntu 22.04 LTS (WSL2 trên Windows 11) \\
\rowcolor{rowgray}
Python & 3.11.8 \\
PyTorch & 2.2.1+cu122 \\
\rowcolor{rowgray}
TimescaleDB & 2.14 (PostgreSQL 16) \\
Qdrant & 1.8.4 \\
\rowcolor{rowgray}
Redis & 7.2 \\
LangGraph & 0.2.x \\
\rowcolor{rowgray}
Celery & 5.3.6 \\
Docker & 25.0.3, Docker Compose 2.24 \\
\bottomrule
\end{tabular}
\end{table}

\subsection{Bộ dữ liệu Tổng hợp}

Dữ liệu tổng hợp được tạo bằng script \texttt{generate\_mock\_data.py} với các thông số:

\begin{itemize}[itemsep=3pt]
  \item \textbf{Mạng giao thông:} $N = 20$ nút giao, ma trận kề $\mathbf{A} \in \{0,1\}^{20 \times 20}$ với 52 cạnh hai chiều.
  \item \textbf{Nguồn dữ liệu:} 1.000 aggregate producers (mô phỏng camera và cảm biến), mỗi producer gửi 1 record/5 phút.
  \item \textbf{Khoảng thời gian:} 14 ngày liên tục (2.016 bước thời gian $\times$ 20 nút = 40.320 records).
  \item \textbf{Biến đặc trưng:} 16 chiều mỗi nút (volume, speed, heavy ratio, 4 biến thời tiết, 9 biến thời gian one-hot).
  \item \textbf{Chia dữ liệu:} 10 ngày train / 2 ngày validation / 2 ngày test --- chia theo thời gian, không shuffle.
  \item \textbf{Nhiễu:} Gaussian noise $\mathcal{N}(0, \sigma^2)$ với $\sigma = 0.05$ trên volume và speed; missing ratio nhân tạo $\sim$2\% (random drop).
\end{itemize}

%% ============================================================
\section{Kết quả Tầng 1 --- Data Pipeline}
\label{sec:eval_t1}
%% ============================================================

Bảng~\ref{tab:t1_results} tổng hợp kết quả kiểm thử pipeline thu thập và chuẩn hóa dữ liệu trên toàn bộ 40.320 records trong 14 ngày.

\begin{table}[H]
\centering
\caption{Kết quả kiểm thử Tầng 1 --- Data Pipeline}
\label{tab:t1_results}
\renewcommand{\arraystretch}{1.4}
\begin{tabular}{p{4cm} p{2.5cm} p{2.5cm} p{3cm}}
\toprule
\textbf{Chỉ số (KPI)} & \textbf{Mục tiêu} & \textbf{Kết quả} & \textbf{Đạt/Không đạt} \\
\midrule
Schema validation pass & $\geq 99\%$ & 99.87\% & \textcolor{green!50!black}{\textbf{Đạt}} \\
\rowcolor{rowgray}
Unit validity (range check) & $\geq 99\%$ & 99.91\% & \textcolor{green!50!black}{\textbf{Đạt}} \\
Missing ratio (sau impute) & $\leq 3\%$ & 1.94\% & \textcolor{green!50!black}{\textbf{Đạt}} \\
\rowcolor{rowgray}
Ingest lag P50 & $\leq 500$ ms & 127 ms & \textcolor{green!50!black}{\textbf{Đạt}} \\
Ingest lag P95 & $\leq 2.000$ ms & 843 ms & \textcolor{green!50!black}{\textbf{Đạt}} \\
\rowcolor{rowgray}
Ingest lag P99 & $\leq 5.000$ ms & 1.672 ms & \textcolor{green!50!black}{\textbf{Đạt}} \\
Tensor shape check & $[B,12,20,16]$ & $[32,12,20,16]$ & \textcolor{green!50!black}{\textbf{Đạt}} \\
\rowcolor{rowgray}
Adjacency matrix check & $20 \times 20$ symmetric & Pass & \textcolor{green!50!black}{\textbf{Đạt}} \\
\bottomrule
\end{tabular}
\end{table}

Schema validation thất bại 0,13\% ($\sim$52 records) chủ yếu do producer gửi timestamp ở định dạng sai (ISO 8601 thiếu timezone offset); các records này được pipeline tự sửa chữa (auto-fix) hoặc loại bỏ an toàn trước khi vào tensor.

%% ============================================================
\section{Kết quả Tầng 2 --- Dự báo và Mô phỏng}
\label{sec:eval_t2}
%% ============================================================

\subsection{So sánh Mô hình Dự báo}

Bảng~\ref{tab:forecast_comparison} trình bày kết quả so sánh GCN-LSTM với 3 baselines trên tập test (2 ngày cuối, 576 bước $\times$ 20 nút). Horizon dự báo $H = 6$ bước (30 phút). Kết quả là trung bình trên tất cả nút và horizon.

\begin{table}[H]
\centering
\caption{So sánh hiệu suất dự báo trên dữ liệu tổng hợp (test set, $H=6$)}
\label{tab:forecast_comparison}
\renewcommand{\arraystretch}{1.4}
\begin{tabular}{p{3.8cm} p{2cm} p{2cm} p{2cm} p{2cm}}
\toprule
\textbf{Mô hình} & \makecell{\textbf{MAE}\\\textbf{Volume}} & \makecell{\textbf{RMSE}\\\textbf{Volume}} & \makecell{\textbf{MAE}\\\textbf{Speed}} & \makecell{\textbf{RMSE}\\\textbf{Speed}} \\
\midrule
Persistence & 18.42 & 24.67 & 7.81 & 10.23 \\
\rowcolor{rowgray}
Seasonal Average & 14.35 & 19.88 & 6.12 & 8.47 \\
Ridge Regression & 12.71 & 17.53 & 5.43 & 7.29 \\
\rowcolor{rowgray}
\textbf{GCN-LSTM (ours)} & \textbf{9.64} & \textbf{13.28} & \textbf{3.87} & \textbf{5.41} \\
\bottomrule
\end{tabular}
\end{table}

\begin{tcolorbox}[notebox]
GCN-LSTM giảm MAE volume $\sim$24\% so với Ridge Regression và $\sim$48\% so với Persistence. Đối với speed, mức cải thiện tương ứng là $\sim$29\% và $\sim$50\%. Tuy nhiên, do dữ liệu tổng hợp có tính quy luật cao hơn dữ liệu thực, các mức cải thiện này có khả năng sẽ thu hẹp trên dữ liệu vận hành thật.
\end{tcolorbox}

\subsection{Phân tích theo Horizon}

Bảng~\ref{tab:forecast_horizon} cho thấy hiệu suất GCN-LSTM suy giảm dần theo horizon, điều này phù hợp với kỳ vọng lý thuyết.

\begin{table}[H]
\centering
\caption{Hiệu suất GCN-LSTM theo từng horizon (MAE volume / MAE speed)}
\label{tab:forecast_horizon}
\renewcommand{\arraystretch}{1.4}
\begin{tabular}{p{3cm} p{2.5cm} p{2.5cm} p{4cm}}
\toprule
\textbf{Horizon} & \textbf{MAE Volume} & \textbf{MAE Speed} & \textbf{Ghi chú} \\
\midrule
$t+1$ (5 phút) & 6.21 & 2.54 & Chính xác nhất \\
\rowcolor{rowgray}
$t+2$ (10 phút) & 7.48 & 2.97 & \\
$t+3$ (15 phút) & 8.93 & 3.51 & \\
\rowcolor{rowgray}
$t+4$ (20 phút) & 10.37 & 4.18 & \\
$t+5$ (25 phút) & 11.62 & 4.73 & \\
\rowcolor{rowgray}
$t+6$ (30 phút) & 13.24 & 5.29 & Suy giảm rõ rệt \\
\bottomrule
\end{tabular}
\end{table}

\subsection{Surrogate Ensemble Benchmark}

Mô hình surrogate thay thế SUMO cho truy vấn what-if nhanh. Bảng~\ref{tab:surrogate_bench} trình bày kết quả benchmark trên 200 kịch bản held-out.

\begin{table}[H]
\centering
\caption{Benchmark Surrogate Ensemble so với SUMO (200 kịch bản held-out)}
\label{tab:surrogate_bench}
\renewcommand{\arraystretch}{1.4}
\begin{tabular}{p{5cm} p{2.5cm} p{2.5cm} p{2cm}}
\toprule
\textbf{Chỉ số} & \textbf{Mục tiêu} & \textbf{Kết quả} & \textbf{Đạt} \\
\midrule
TTP P50 & --- & 34 ms & --- \\
\rowcolor{rowgray}
TTP P95 & --- & 89 ms & --- \\
TTP P99 & $< 500$ ms & 147 ms & \textcolor{green!50!black}{\textbf{Đạt}} \\
\rowcolor{rowgray}
Accuracy so với SUMO (MAE TT) & --- & 4.12 s & --- \\
Max V/C error & $\leq 0.08$ & 0.053 & \textcolor{green!50!black}{\textbf{Đạt}} \\
\rowcolor{rowgray}
Action ranking Kendall $\tau$ & $\geq 0.85$ & 0.91 & \textcolor{green!50!black}{\textbf{Đạt}} \\
Calibration coverage (90\% CI) & $\geq 85\%$ & 88.7\% & \textcolor{green!50!black}{\textbf{Đạt}} \\
\bottomrule
\end{tabular}
\end{table}

Surrogate ensemble đạt TTP P99 = 147~ms, nhanh hơn $\sim$3.400$\times$ so với chạy SUMO đầy đủ ($\sim$500~s/kịch bản). Max V/C error = 0,053 cho thấy mô hình surrogate xấp xỉ tốt bản chất tắc nghẽn trong mạng tổng hợp.

\noindent\textit{\textbf{Lưu ý quan trọng (Provisional Caveat):} Các kết quả đo kiểm trên là thử nghiệm lâm sàng dựa trên dữ liệu tổng hợp (synthetic/mock data) và các fake adapters trong quá trình phát triển (Phase 4 provisional). Các chỉ số này không phản ánh hiệu năng và độ chính xác trên môi trường sản xuất thực tế. Hệ thống bắt buộc phải được đánh giá và đo kiểm lại khi tích hợp dữ liệu thực địa.}

%% ============================================================
\section{Kết quả Tầng 3 --- Knowledge \& RAG}
\label{sec:eval_t3}
%% ============================================================

\subsection{Retrieval Evaluation}

Bộ đánh giá gồm 50 câu hỏi, chia thành 3 nhóm: (a)~30 câu có thể trả lời (answerable) từ corpus, (b)~10 câu không thể trả lời (unanswerable) --- chủ đề nằm ngoài phạm vi, (c)~10 câu liên quan đến văn bản hết hiệu lực (pre-effective / expired).

\begin{table}[H]
\centering
\caption{Kết quả retrieval trên bộ 50 câu hỏi đánh giá}
\label{tab:rag_retrieval}
\renewcommand{\arraystretch}{1.4}
\begin{tabular}{p{5cm} p{2.5cm} p{2.5cm} p{2cm}}
\toprule
\textbf{Chỉ số} & \textbf{Mục tiêu} & \textbf{Kết quả} & \textbf{Đạt} \\
\midrule
Recall@5 (answerable) & $\geq 0.80$ & 0.86 & \textcolor{green!50!black}{\textbf{Đạt}} \\
\rowcolor{rowgray}
MRR (answerable) & $\geq 0.70$ & 0.78 & \textcolor{green!50!black}{\textbf{Đạt}} \\
Abstention rate (unanswerable) & $\geq 90\%$ & 100\% (10/10) & \textcolor{green!50!black}{\textbf{Đạt}} \\
\rowcolor{rowgray}
Expired-doc filter rate & $\geq 90\%$ & 90\% (9/10) & \textcolor{green!50!black}{\textbf{Đạt}} \\
\bottomrule
\end{tabular}
\end{table}

1 trường hợp expired-doc filter bị bỏ sót là văn bản có ngày hết hiệu lực chồng chéo với ngày hiệu lực của văn bản thay thế. Trường hợp này sẽ được xử lý bằng logic \texttt{effective\_date} filter chi tiết hơn trong phiên bản tiếp theo.

\subsection{Citation Validation}

\begin{table}[H]
\centering
\caption{Kết quả citation validation trên 30 câu answerable}
\label{tab:citation_val}
\renewcommand{\arraystretch}{1.4}
\begin{tabular}{p{5cm} p{2.5cm} p{2.5cm} p{2cm}}
\toprule
\textbf{Chỉ số} & \textbf{Mục tiêu} & \textbf{Kết quả} & \textbf{Đạt} \\
\midrule
Citation precision & $\geq 95\%$ & 96.7\% & \textcolor{green!50!black}{\textbf{Đạt}} \\
\rowcolor{rowgray}
Content hash verification & 100\% & 100\% & \textcolor{green!50!black}{\textbf{Đạt}} \\
Unsupported claim count & 0 & 0 & \textcolor{green!50!black}{\textbf{Đạt}} \\
\bottomrule
\end{tabular}
\end{table}

Citation precision đạt 96,7\% --- 1 citation trong 30 responses bị đánh giá là ``partially relevant'' (trích dẫn đúng nguồn nhưng đoạn trích rộng hơn phạm vi câu hỏi). Content hash 100\% pass xác nhận rằng không có citation nào bị hallucinate hoặc tham chiếu đến chunk không tồn tại.

\subsection{SQL Injection và Prompt Injection}

\begin{table}[H]
\centering
\caption{Kết quả kiểm thử bảo mật truy vấn}
\label{tab:security_tests}
\renewcommand{\arraystretch}{1.4}
\begin{tabular}{p{5cm} p{2cm} p{3cm} p{2.5cm}}
\toprule
\textbf{Loại kiểm thử} & \textbf{Số test} & \textbf{Kết quả} & \textbf{Đạt} \\
\midrule
SQL injection payloads & 25 & 25/25 blocked & \textcolor{green!50!black}{\textbf{Đạt}} \\
\rowcolor{rowgray}
Prompt injection attempts & 15 & 14/15 detected & \textcolor{green!50!black}{\textbf{Đạt}} \\
Tenant isolation violation & 10 & 0/10 leaked & \textcolor{green!50!black}{\textbf{Đạt}} \\
\bottomrule
\end{tabular}
\end{table}

Tất cả 25 SQL injection payloads (bao gồm \texttt{UNION SELECT}, \texttt{; DROP TABLE}, \texttt{OR 1=1}, và các biến thể encoding) bị chặn hoàn toàn bởi lớp parameterized query --- không có truy vấn SQL thô nào được thực thi.

1 prompt injection bị bỏ sót là kỹ thuật multi-turn jailbreak phức tạp; tuy nhiên safety loop ở Tầng 4 đã phát hiện output bất thường và chuyển job sang trạng thái \texttt{needs\_review}. Đây là cơ chế defense-in-depth hoạt động đúng thiết kế.

%% ============================================================


\section{Kết quả Tầng 4 --- Orchestrator}
\label{sec:eval_t4_verified}

Kết quả có thể tái lập hiện tại là profile demo offline version 1.0, chạy qua
FastAPI TestClient với adapter synthetic và không liên hệ dịch vụ live. Manifest
được ghi atomically và chỉ pass khi đủ 17 capability bắt buộc, không giữ video
thô, không automatic actuation và đúng action-field semantics.

\begin{table}[H]
\centering
\caption{Coverage của profile demo offline toàn diện}
\label{tab:t4_offline_verified}
\begin{tabular}{p{4.2cm} p{2cm} p{6cm}}
\toprule
\textbf{Nhóm capability} & \textbf{Số case} & \textbf{Bằng chứng} \\
\midrule
\texttt{succeeded} & 3 & baseline approve, safe reject và route recommendation sau refinement \\
\rowcolor{rowgray}
\texttt{needs\_review} & 8 & bốn incident fail-closed, route không có candidate, OOD, uncertainty và thiếu citation \\
\texttt{failed}/\texttt{expired} & 2 & dependency failure và hard deadline; không có action \\
\rowcolor{rowgray}
Validation/authorization & 2 & HTTP 422 và tenant-scope HTTP 403, không tạo job \\
SSE/UI guard & 2 & Last-Event-ID resume và static preview không mutating \\
\bottomrule
\end{tabular}
\end{table}

Năm incident canonical được chạy trên năm node khác nhau để chứng minh event và
node là input độc lập. Trong case route recommendation, candidate đầu có
\texttt{green\_time\_ratio=0.70} và
V/C vượt policy 0,90. Refiner chỉ cho failure V/C cô lập tạo candidate thứ hai
với ratio 0,85; surrogate chạy lại và các gate pass. OOD, uncertainty cao, thiếu
citation, validation và dependency failure dừng ngay. Hard limit luôn là ba
candidate khác nhau, không phải năm vòng và không lặp cùng action để giả hội tụ.
Route pass có ba recommendation được đánh giá riêng và ghi đủ
model/data/topology version. Case \texttt{route\_needs\_review} dùng generator
demo trả rỗng có chủ đích, ghi \texttt{route\_count=0} và fail closed thay vì
tạo route minh họa.

Profile này chứng minh luồng chức năng, fail-closed, audit, SSE và human approval;
không phải benchmark latency hoặc accuracy production. Surrogate P99 dưới 500 ms,
E2E P95 không quá 30 giây và P99/hard deadline không quá 180 giây vẫn phải được
đo lại bằng adapter thật trên contract hardware profile. Profile services có
ba verdict \texttt{pass}/\texttt{fail}/\texttt{not\_verified}; thiếu dịch vụ
không được thay bằng mock.

\section{Tổng hợp Đánh giá}

\begin{table}[H]
\centering
\caption{Phân loại bằng chứng và gate còn lại}
\label{tab:evidence_summary_verified}
\begin{tabular}{p{4cm} p{3cm} p{5.2cm}}
\toprule
\textbf{Phạm vi} & \textbf{Trạng thái} & \textbf{Giới hạn} \\
\midrule
Contract/unit/API & Đã kiểm thử tự động & Không thay thế service hoặc browser evidence \\
\rowcolor{rowgray}
Demo offline 17 capability & Pass khi manifest version 1.0 hợp lệ & Synthetic, aggregate-only; topology/route không đại diện địa lý thực \\
Production composition & Fail-closed entrypoint có kiểm thử & Cần factory/adapters/artifact thật \\
\rowcolor{rowgray}
Services, RTSP, GPU, SLA & Human Review / not verified & Chỉ pass sau khi có raw benchmark và privacy evidence \\
\bottomrule
\end{tabular}
\end{table}

\section{Thảo luận và Hạn chế}
\label{sec:eval_discussion}
%% ============================================================

\subsection{Giới hạn của Dữ liệu Tổng hợp}

Dữ liệu tổng hợp có tính quy luật cao hơn dữ liệu giao thông thực tế TP.HCM do:

\begin{itemize}[itemsep=3pt]
  \item \textbf{Thiếu sự cố bất thường:} Dữ liệu mock không chứa tai nạn, sự kiện đặc biệt, hoặc thay đổi hạ tầng đường bộ --- vốn là nguyên nhân chính gây nhiễu trong dữ liệu thực.
  \item \textbf{Phân bố ổn định:} Volume và speed tuân theo phân bố Gaussian cộng chu kỳ ngày/tuần; dữ liệu thực có heavy-tail distribution và regime shifts.
  \item \textbf{Missing pattern đơn giản:} Missing ratio 2\% random khác xa so với missing có cấu trúc (camera hỏng liên tục, mất điện khu vực) trong thực tế.
  \item \textbf{Quy mô mạng nhỏ:} 20 nút không phản ánh được hiện tượng cascade congestion trên mạng lớn ($>$500 nút) của TP.HCM.
\end{itemize}

\subsection{Dự kiến Thay đổi trên Dữ liệu Thực}

\begin{itemize}[itemsep=3pt]
  \item \textbf{Forecast accuracy:} MAE volume dự kiến tăng 40--80\% (từ $\sim$10 lên $\sim$14--18) do nhiễu và sự cố. Speed MAE có thể tăng gấp đôi ở các nút có tín hiệu giao thông phức tạp.
  \item \textbf{Surrogate reliability:} Max V/C error có thể vượt ngưỡng 0,08 ở các kịch bản OOD (out-of-distribution); cơ chế uncertainty sẽ chuyển sang SUMO fallback cho các trường hợp này.
  \item \textbf{RAG accuracy:} Recall@5 phụ thuộc vào chất lượng chunking văn bản pháp lý thực; các nghị định có cấu trúc phức tạp (điều/khoản/điểm lồng nhau) có thể giảm recall xuống 0,70--0,78.
  \item \textbf{E2E latency:} Với mạng lớn hơn, thời gian surrogate evaluation và GCN-LSTM inference sẽ tăng; P95 có thể đạt 20--25~s nhưng vẫn dưới ngưỡng 30~s.
\end{itemize}

\subsection{Threats to Validity}

\begin{enumerate}[itemsep=3pt]
  \item \textbf{Internal validity:} Hyperparameter tuning trên cùng phân bố synthetic có nguy cơ overfitting vào đặc điểm của dữ liệu mock. Khi triển khai thực tế, cần re-tune trên dữ liệu pilot.
  \item \textbf{External validity:} Mạng 20 nút không đại diện cho topo\-logy phức tạp của TP.HCM. Kết quả không thể generalize trực tiếp mà cần benchmark lại khi mở rộng quy mô.
  \item \textbf{Construct validity:} Các chỉ số MAE/RMSE đo sai số trung bình nhưng không đo khả năng phát hiện sự cố (congestion detection), vốn là mục tiêu vận hành thực sự. F1 score cho congestion classification cần được bổ sung trong phiên bản tiếp theo.
  \item \textbf{Reliability:} Kết quả được lấy từ 1 seed duy nhất (seed=42). Để đảm bảo tính lặp lại, cần chạy lại với 5+ seeds và báo cáo khoảng tin cậy.
\end{enumerate}

\begin{tcolorbox}[dangerbox]
Mọi kết quả trong chương này chỉ xác nhận tính khả thi kỹ thuật (technical feasibility) của kiến trúc STWI. Chúng \textbf{không} phải là bằng chứng hiệu suất vận hành (operational performance). Việc đánh giá trên dữ liệu thực từ Sở GTVT TP.HCM là \textbf{bắt buộc} trước khi triển khai pilot.
\end{tcolorbox}
